\documentclass[11pt]{article}
\usepackage[letterpaper,margin=1in]{geometry}
\usepackage{amsmath,graphicx}
\usepackage[colorlinks=true,urlcolor=blue,linkcolor=blue]{hyperref}
\setlength{\parindent}{0pt}
\setlength{\parskip}{0.65em}
\pagestyle{plain}
\title{Finding hidden groups in a network\\A plain-language guide to the research manuscript}
\author{Companion document}
\date{November 15, 2026}
\begin{document}
\maketitle

\section*{1. What problem does the manuscript study?}
A network is a collection of points joined by links. The points might
represent people, computers, or other objects. A link records a
connection between two of them. Mathematicians call the points
\emph{vertices}, the links \emph{edges}, and the whole network a
\emph{graph}. A point's \emph{degree} is its number of neighbors.

The problem is to find groups whose labels have been erased. The
original hidden grouping is called the \emph{plant}. The
manuscript studies a particularly precise kind of hidden group: every
member of one group has exactly the same number of neighbors in each
group. These numbers are exact, rather than approximate averages.
The actual neighbors can differ from one member to another.

For example, suppose every point has three neighbors: exactly two in
its own group and one in the other group. The task is to recover the
two groups from the links alone. A small local picture is:

\begin{center}
\setlength{\unitlength}{1mm}
\begin{picture}(145,44)
\put(3,4){\framebox(62,35){}}
\put(81,4){\framebox(61,35){}}
\put(6,34){\makebox(55,3){Group A}}
\put(84,34){\makebox(55,3){Group B}}
\put(40,21){\circle*{2}}
\put(18,10){\circle*{2}}
\put(18,32){\circle*{2}}
\put(106,21){\circle*{2}}
\put(18,10){\line(2,1){22}}
\put(18,32){\line(2,-1){22}}
\put(40,21){\line(1,0){66}}
\put(29,9){\makebox(29,3){one point}}
\end{picture}

\small Only the three links of one point are shown; the other links
and points have been omitted.
\end{center}

The result concerns networks generated by a specified random model.
It gives a mathematical guarantee for that model, rather than a test
that arbitrary social or computer networks must pass. Its central
achievement is to combine exact group recovery with a proved bound
on the algorithm's average work, even when the network stays sparse.
Here \emph{sparse} means that each point has only a fixed number of
neighbors as the network grows.

\newpage
\section*{2. Exactly which networks does the guarantee cover?}
The model starts with a fixed number of hidden groups and a table of
neighbor counts. A row of the table says how many neighbors a member
of that group has in each group. The manuscript calls this table a
\emph{degree frame}. The algorithm is not given the table, the group
sizes, or the actual number of groups. It receives fixed upper bounds
on the number of groups and on the number of neighbors of a point.

The sizes and table must fit together. For example, the number of links
counted from group A to group B must equal the number counted back
from B to A. The groups may have different sizes, and different groups
may have different total numbers of neighbors. Their size ratios stay
fixed as the network grows. Links join distinct points, and there is
at most one link between a given pair.

The table must connect the groups: one can move between any two
groups through a sequence of cross-group connection types with
positive counts. The one-group case is allowed. This condition is
about the table; it does not require every sampled network to be connected.

Within each group, the internal links are chosen uniformly among all
possibilities with the required degrees. Links between each pair of
groups are chosen uniformly subject to their required counts. These
different choices are independent. Those randomness assumptions are
part of the theorem.

There is also a condition on the count table. In plain terms, one must
be able to give each type of connection a weight, the same when viewed
from either end. Each weight is at least zero and strictly less than
one half. At every point, the weights of all its incident links must
add to one. The weights only certify that the theorem applies; the
algorithm does not need them as input.

An important broad case is easy to recognize: if every point has the
same fixed degree $d\ge3$, give every connection weight $1/d$.
The condition then holds, including when there are no internal links
within a group. The manuscript also covers some unequal-degree cases.
It does not assert that minimum degree three alone is sufficient.

Another useful example is a \emph{random lift}. Start with a small
template network. Replace each template point by $N$ points. Replace
each template link by a random one-to-one pairing between its two
groups of copies. These groups of copies are called \emph{fibers}.
The theorem covers
every fixed connected simple
template whose points all have the same degree at least three.
The word \emph{fixed} matters: the template and the allowed group
bound stay fixed while $N$ grows. Every admissible size is included;
the guarantee is not limited to one parity of the group size.

\newpage
\section*{3. How can the algorithm find groups without their table?}
The neighbor-count rule creates exact equations. To see the idea,
imagine marking the points in a proposed group with a $1$ and all
other points with a $0$. Adding the marks of a point's neighbors
counts how many of them belong to that proposed group. A correct
partition makes those counts agree with its proposed table.

There are only finitely many possible tables when the maximum degree
and maximum number of groups are fixed. The algorithm can therefore
try all of them. That finite collection can be large; its size is
absorbed into constants depending on the fixed bounds.

The difficult part is searching through the possible assignments of
points to groups. Trying every assignment separately would take far
too long. Instead, the algorithm first solves a simpler version of
the equations using arithmetic on a small repeating clock. For
example, on a three-position clock, adding one to two gives zero.
Mathematicians call this arithmetic over a \emph{finite field}.

This step describes a potentially huge set of assignments through a
smaller collection of basic solutions. Every genuine partition is
represented. The clock arithmetic can also introduce false candidates,
so it is a way to organize the search, rather than the final answer.
The algorithm certifies that its collection of basic solutions is
complete, generates the represented candidates, and checks their
original, ordinary integer neighbor counts. A false modular solution
cannot pass that final check.

The proof bounds the average number of candidates in the compressed
search. It includes rare inputs that have unusually many solutions;
their work is counted in the expectation. Complete certification is
important here. Finding one plausible solution, or assuming that the
solution space is small, would not establish the stated result.

The algorithm lists every valid partition with at most the supplied
number of groups, denoted $K$ in the paper. It then keeps one with
the largest number of groups. With probability tending to one in the
model, every valid partition on this list comes from merging whole
planted groups. Consequently the largest one is precisely the hidden
plant, apart from the arbitrary names of its groups.

This is what the paper means by the \emph{finest bounded-class}
partition. The bound is essential: if one allowed as many groups as
points, putting every point in a separate group would always satisfy
an exact neighbor-count description.

\newpage
\section*{4. What do ``nearly quadratic'' and ``gap-free'' mean?}
Let $n$ be the total number of points. The algorithm's expected work
is nearly proportional to $n^2$, allowing additional factors involving
logarithms of $n$. Thus, for fixed model parameters and ignoring those
logarithmic factors, doubling the network size means roughly four
times as much work in the bound. The expected working memory is
nearly proportional to $n$.

This is an asymptotic mathematical bound, not a measured running-time
prediction for a particular implementation. Constants can depend on
the degree bound, group bound, and other fixed parameters. The paper
does not show that those constants are small.

The word \emph{expected} also matters. Some inputs may have many valid
answers or a much larger intermediate search. The theorem does not
give a quadratic worst-case time limit on every input. It does give
an always-correct randomized listing procedure: randomness affects
how long its certified computations take. It does not permit an
unchecked candidate to be reported as a valid partition. A separate
version imposes a deterministic polynomial work limit, at the cost of
a small probability of interruption under the random model.

Correct listing and correct hidden-group recovery are different.
Every listed answer satisfies the exact equations. On a rare graph,
several answers can satisfy them, and the observed graph alone may
not reveal which one was planted. That is an information limitation,
even for an algorithm with unlimited computing power.

Why avoid an eigenvalue-gap assumption? Many familiar network methods
use special numbers and directions associated with the adjacency
matrix, called eigenvalues and eigenvectors. Their guarantees often
need a signal to be separated from the rest of the spectrum. That
separation can fail in sparse models.

The recovery algorithm here uses exact neighbor-count equations and
integer verification. Its guarantee does not require a numerical gap
around a planted eigenvalue. Some separate results in the manuscript
also show that inherited template eigenvalues acquire no additional
copies in a random regular lift, with high probability. This is a
statement about exact equality of eigenvalues; it does not give a fixed
lower bound on the distance between nearby ones. Neither result
establishes stability when the network is perturbed.

\newpage
\section*{5. Some ambiguity is unavoidable, and its size can be computed}
Return to the example with two equally sized groups, two neighbors
inside one's own group, and one across. Let $N$ be the size of
\emph{one} group, so the full network has $2N$ points. Choose the two
equal-sized hidden groups uniformly from all possibilities before
revealing the network.

The manuscript proves that, as $N$ grows, the probability that another
valid unordered two-group partition exists is asymptotic to
\[
 \frac{4}{N^2}.
\]
Here \emph{asymptotic to} means that the ratio of the true probability
to this expression tends to one. It is not an exact formula at a
particular finite size.

Most ambiguous graphs at this leading order have exactly two possible
partitions. Under this random way of hiding the groups, these two answers
are equally likely after the graph is observed. On those ambiguous
graphs, no estimator can select the planted answer with probability
better than one half. The smallest possible average
recovery error is therefore asymptotic to
\[
 \frac{2}{N^2}.
\]
The listing algorithm, followed by selection of any verified two-group
partition with these neighbor counts,
attains this optimal average error for this way of hiding the groups. The manuscript
computes the corresponding leading constants for all fixed two-group
degree choices with at least one neighbor across and at least three
neighbors in total, including
choices with more links across groups than within them.

The boundary cases explain why the assumptions matter. A cycle is a
closed ring of points. If the known template is a cycle, an unconditioned
random lift is a collection of cycles. Each component can carry its own shift and direction of the
hidden group labels. As the lift grows, even the best possible
hidden-group recovery success tends to zero. If one instead conditions
the lift to be connected and supplies the cycle template, the unordered
partition is uniquely determined. These are different experiments.

There are also fixed unequal-degree templates of minimum degree three
with much stronger ambiguity. The template $K_{3,12}$ has three points
on one side and twelve on the other, with every cross-side pair linked
and no links on a side. Even revealing the three fibers corresponding
to the template's three-point side leaves
exponentially many possibilities for the other fibers, typically.
When the remaining hidden partition is chosen uniformly, the optimal
recovery success is exponentially small. This is why
the positive theorem uses a condition on the whole count table, rather
than treating minimum degree as a sufficient condition by itself.

\newpage
\section*{6. What would make this result important?}
The potential contribution is a rigorous guarantee across a broad
fixed-parameter class of sparse exact-count networks. It joins three
questions that need separate answers: whether the hidden groups are
usually identifiable, whether an algorithm can list all valid answers,
and whether its work can be bounded after charging even the unusual
inputs. It also identifies specific models where recovery is impossible.

The regular-template scope is particularly concrete. The same
guarantee covers every fixed connected simple regular template of
degree at least three, rather than only a few specially chosen small
templates. The algorithm need not be supplied with the template.
For these noiseless models it can work without requiring the planted
spectral signal to stand outside a bulk of other eigenvalues.

These claims should be read alongside their ancestors. Fellows and
Koblitz discussed hidden-fiber recovery for a four-point template with
every pair joined in their 1993 \emph{Kid Krypto} work. A 2023 evaluation by
Bernstein, Cottaar and Lange used Gaussian elimination and grouping
equal algebraic expressions to expose those fibers. These are direct
predecessors of the algebraic recovery idea. The manuscript also uses
classical graph counting, entropy, finite-field moments, local limit
estimates, and certified linear algebra. It does not claim to have
invented these tools.

The possible advance lies in the exact scope and the proved composition
of the ingredients: unknown tables, bounded numbers of groups,
independent sparse factors, all admissible sizes, complete certified
listing, expected work, and sharp ambiguity constants. Establishing
historical firstness would require a further literature assessment.
The manuscript explicitly does not certify it.

There are clear limits. The number of groups and degree bounds do not
grow with network size. The randomness is specified and independent.
Neighbor counts are exact. The results do not promise recovery after
missing links, added links, measurement errors, or other noise. A
network that approximately resembles the model is not automatically
covered. The paper also makes no claim of practical benchmark speed,
an arbitrary-network clustering method, or a general classification
of all recoverable templates.

This is a research manuscript supported by internal mathematical
audits and finite checks. Those checks help find mistakes; they are
not formal verification or external peer review. This companion is
an explanation of the stated guarantees and their limits. The main
manuscript contains the precise hypotheses, proofs, and references.

\end{document}
